\scanpage{162}
modules dans les $\underline U$ modules, i.e. un
\deleted{\unsure{isomph|isomorphisme}} homomorphisme $\underline U^\circ\longrightarrow\underline U$, i.e.
un \deleted{\unsure{anti-automorphisme}}\inserted{application $\sigma\colon\underline U\longrightarrow\underline U$ qui est \unsure{$k$-linéaire|$k$-lin.}}
et anti-multiplicative; $\breve F\breve F=I$ signifie
que $\sigma^{\unsure{2}}=1$, donc c'est un \uline{antiautomorphisme involutif}, \uline{noté $x\longrightarrow\breve x$}. Réciproquement,
un tel antiautomorphisme définit
\deleted{\unsure{une}} un foncteur $K^{\underline U}\longrightarrow K^{\underline U}$. Dans l'hyp. ci-dessus, \deleted{\unsure{si|qui}} --- Si on a de
plus une augmentation $\varepsilon$, la ``\uline{compatibilité}
\uline{de $\sigma$ avec $\varepsilon$}, i.e. $\varepsilon(x)=\varepsilon(\breve x)$
signifie que le \deleted{$\underline U$-structure de}\inserted{\uline{dual du $\underline U$ module}}
\deleted{\unsure{du $k$-fonct}}\uline{trivial est le $U$-module trivial}. \deleted{En} Si
on a un \uline{hom} continu $\varphi:\underline V\longrightarrow\underline U$, le
fait qu'il est \uline{compatible avec $\sigma$},
\deleted{Supposons donné un bifoncteur}
i.e. $\varphi(\breve x)=\overset{\vee}{\overline{\varphi(x)}}$, signifie que
\deleted{$F(A,B)$}
\deleted{si} \deleted{$\overline\varphi(V)$} si $V\in K^{\underline U}$, on a
$\overline\varphi(V)'=\overline\varphi(V')$, i.e le foncteur
permis est compatible avec le
passage au dual.

% archive sheet number: 161
